\scanpage{24}
lieu d'expliciter

a) la correspondance entre \inserted{les $8$} sommets du cube, et les \inserted{$8$} éléments d'ordre $3$ de $G$ (``rotations-sommets'')

b) \deleted{les éléments}\inserted{la corr.} entre les $6$ paires d'arêtes antipodiques du cube, et les $6$ \deleted{\unsure{réflexions}} réflexions de $G=\unsure{\mathcal{G}_I|\mathfrak{S}_I}$

c) la corr. entre les $6$ \deleted{éléments de $\operatorname{Circ}(I)$, et}\inserted{faces du cube, et} les \deleted{$6$ \unsure{deux éléments de $G$}}\inserted{permutations circulaires} (él. d'ordre $4$ de $G\simeq\operatorname{Aut}^+(\text{Cube})$) correspondants,

et de décrire également une \uline{orientation canonique} du cube, ou ce qui revient au $=$, une \deleted{trivialisation} iso.
\[
\underset{\substack{\Vert\mathrm{def}\\\underline{\omega}_J(\simeq\underline{\omega}_I)}}
{\underline{\omega}_{\unsure{C|Q}}}
\simeq \bigwedge_{i\in J}
\underset{\substack{\text{ens des deux \deleted{\unsure{éléments}}\inserted{autom.} d'ordre $4$}\\
\text{dans la structure carré}\\
Q_i\text{ définie par }i\in J\\
\text{i.e. ens des deux orientations}\\
\text{de }Q_i.}}{\underline{\omega}_i}
\]

Comme on a
\[
\underline{\omega}_i\simeq
\underset{\substack{\text{ens. des deux}\\\text{codiagonales de }Q_i}}
{\underline{\omega}}\wedge
\underset{\substack{\text{ens des diag.}\\
\text{pour le carré }Q_i,\ \text{i.e. }D_i=
I/\underset{\text{translation indiquée par }i}{\underline{a}_i}}}{D_i}
\]
la formule \deleted{\unsure{antérieure}} s'écrit aussi
\[
\underset{\substack{\text{\unsure{torseur}}\\\text{sous }\mathbb{F}_2^J}}
{\bigwedge_{i\in J}D_i}\simeq\underline{1}
\]
Or on a \deleted{la suite exacte} l'inclusion
\diagram{\begin{array}{ccc}
I&\hookrightarrow&\prod_{i\in J}D_i\\
\substack{\text{torseur}\\\text{sous }V_J(\mathbb{F}_2)}
&\subset&
\substack{\text{torseur sous}\\\mathbb{F}_2^J}
\end{array}}
% archive foliation: 23
